\begin{table}[t]\centering\footnotesize
\caption{Operational rules for the nine reuse bindings. Apply each item to a specified result record and deployment claim.}\label{tab:audit-rules}
\begin{tabularx}{\columnwidth}{@{}p{.30\columnwidth}Y@{}}\toprule
Item & Evidence required for a Bound (\textbf{B}) judgment\\\midrule
R1 Target & Hidden property, inference goal and authorized/protected interest are explicit.\\
R2 Receiver & Observer permission, released representation and auxiliary inputs identify the evaluated view.\\
R3 Intervention & Placement and execution mode bind the transform to stored data, wire emission or executed actions.\\
R4 Protected fit & Attacker fitting/calibration uses the protected condition; training and protected testing are identified.\\
R5 Matched adaptation & Clean- and protected-condition fits are evaluated on the same defended test population.\\
R6 Workload & Tasks, system/generator configuration and holdout specify the supported transfer axis.\\
R7 Inference & Metric, denominator and evaluation condition bind the attack outcome to the scoped mechanism.\\
R8 Utility & Task/component performance is measured at the same protected operating point, with its state boundary explicit.\\
R9 Cost & Resource units, baseline and measurement object price that operating point.\\\bottomrule
\end{tabularx}
\par\smallskip\raggedright\footnotesize \textbf{B}: bound; \textbf{P}: partial binding; \NRR{}: unreported in inspected scope; \NRA{}: source unavailable; \NRV{}: binding unverified; --: inapplicable. First determine applicability, then assign evidence status. Family \textbf{B} means at least one scoped record satisfies the item, not that every result does.
\end{table}
